This chapter reviews the current literature on the four selected optimization techniques, focusing on the performance benefits each method provides in relation to the research questions. Relevant publications were identified on Google Scholar using keyword combinations of the optimization method and ``RAG'' or ``retrieval-augmented generation''. In addition, snowballing \cite{guidelines}, starting from comprehensive surveys (e.g., \cite{gan2025rag, gao2023retrieval, wu2025ragnlp}), was used to discover further relevant studies. Due to the rapid pace of development in this research area, the selected literature consists primarily of articles published after 2023. The remainder of this chapter is structured as follows: Section~\ref{lit:cascading} reviews model cascading and routing, Section~\ref{lit:retrieval} examines small-to-big retrieval, Section~\ref{lit:re_rank} covers semantic re-ranking, and Section~\ref{lit:caching} discusses prompt caching.

\subsection{Cascading}
\label{lit:cascading}
Modern RAG systems encounter diverse workloads, ranging from simple informational queries to highly complex analytical tasks. In production-grade enterprise environments — where high availability, low latency, and cost efficiency are critical constraints — a "one-size-fits-all" architecture is fundamentally sub-optimal \cite{jeong2024adaptive, chen2023frugalgpt}. It over-allocates expensive computational resources for trivial questions, yet may lack the retrieval depth and reasoning capacity required for complex queries. To mitigate this, model cascading or query routing can be employed to dynamically direct inputs to an appropriately scaled configuration of the RAG pipeline \cite{ong2025routellm}. For instance, a simple query can bypass the full pipeline and utilize a lightweight LLM with minimal context. Conversely, a complex query requiring synthesis across multiple enterprise knowledge bases can trigger an augmented pipeline with extended retrieval top-k parameters and a more capable foundation model.

Formally, model cascading and query routing can be structured in two primary ways: classification-based routing and sequential cascading. In a classification-based query routing system, let $\mathcal{M} = \{M_1, M_2, \ldots, M_m\}$ be a set of generator models ordered by capacity and cost. A routing classifier $f_{\text{route}}(q) \in \{1, \ldots, m\}$ maps the query $q$ to a specific model index, yielding:
\begin{equation}
a = M_{f_{\text{route}}(q)}(q, R_k(q, D))
\end{equation}
Alternatively, in a sequential model cascade (e.g., FrugalGPT), queries are evaluated sequentially. Let $S(M_i(q, C))$ be a confidence or scoring function evaluating the quality of the generated output from model $M_i$. The final answer is generated by the first model in the sequence that meets a reliability threshold $\tau_i$:
\begin{equation}
a = M_i(q, R_k(q, D)) \quad \text{where } i = \min \{ j \mid S(M_j(q, R_k(q, D))) \ge \tau_j \}
\end{equation}

Model cascading offers substantial operational cost reductions both in isolated Large Language Model (LLM) invocations and in retrieval-augmented generation (RAG) pipelines whose overall cost profiles are dominated by API usage. For instance, FrugalGPT \cite{chen2023frugalgpt} routes individual requests sequentially from lightweight to heavyweight models based on a reliability score generated by a scoring classifier. This cascading approach achieves up to a 98.3\% cost reduction compared to a homogeneous GPT-4 baseline while maintaining comparable response quality. Similarly, RouteLLM \cite{ong2025routellm} yields a 1.14$\times$ to 3.66$\times$ cost reduction compared to a static GPT-4 deployment by dynamically bypassing expensive models, preserving approximately 95\% of the baseline's quality. Similar cost reduction patterns have also been demonstrated by other adaptive RAG frameworks \cite{jeong2024adaptive, guo2026deepsieve}.

Routing low-complexity queries to smaller models yields significant latency benefits due to their faster inference speeds. Crucially, the computational overhead introduced by the routing decision itself remains minimal, typically amounting to approximately 11\, ms per query \cite{zhang2025ragrouter} or a 0.4\% latency increase \cite{ong2025routellm}. In straightforward scenarios, bypassing retrieval operations entirely can yield order-of-magnitude improvements: \citet{jeong2024adaptive} reports a reduction in generation latency from 3.08 and 27.18 seconds (for single-step and multi-step retrieval pipelines) to 0.35 seconds. Furthermore, adaptive stopping mechanisms, such as that proposed by \citet{park2025stoprag}, can prune unnecessary retrieval-generation loops, reducing the total number of iterations by 14--49\%. 

While providing both operational cost and latency benefits, query routing has a modest impact on quality. For example, \citet{jeong2024adaptive} found that the F1 score remained practically unchanged, while achieving a substantial reduction in latency (~56\%) and in token consumption (63\%) compared to a multi-step retrieval baseline (using GPT-3.5). Similar results have been shown by \citet{guo2026deepsieve}, where token consumption decreased by up to ~90\% while losing only 0.8 points on F1, but improving the Exact Match score by 2.6 points compared to the agentic baseline, and yielding 20.7 and 15.5 percentage point improvements in F1 and Exact Match (EM) scores, respectively, compared to naive RAG.

\subsection{Small-to-Big retrieval}
\label{lit:retrieval}

Document chunking is a foundational phase in RAG pipelines, directly determining the granularity of the vector index. However, selecting a static chunk size introduces a fundamental trade-off between retrieval precision and contextual completeness. While smaller chunks are computationally efficient and exhibit high intrachunk cohesion (ICC), they frequently suffer from low document contextual coherence (DCC) due to semantic fragmentation \cite{demourajunior2026adaptive}. Conversely, larger chunks preserve rich contextual relationships but degrade retrieval accuracy by diluting query-relevant semantic features within dense vector representations. Small-to-big retrieval resolves this tension by decoupling the representation used for vector search from the context provided to the generator. In this architecture, dense retrieval is executed over fine-grained index nodes, while the generation phase utilizes a broader, mapped context (such as parent chunks or entire source documents) retrieved via sibling pointer mapping. Implementations of this technique range from neighboring chunk expansion, which dynamically retrieves adjacent sliding windows to restore structural continuity \cite{lumer2025rethinking}, to hierarchical abstract trees, which map granular leaf nodes to high-level thematic summaries\cite{li2026bridgerag}.

Formally, small-to-big retrieval can be modeled by representing the search space as a set of fine-grained index nodes $D_{\text{small}} = \{d_1, \ldots, d_N\}$ and a context-mapping function $g: D_{\text{small}} \to D_{\text{parent}}$ that maps each small chunk to its corresponding parent chunk or expanded context window. The retriever first operates on the small chunks:
\begin{equation}
C_{\text{small}} = R_k(q, D_{\text{small}})
\end{equation}
The final answer is then generated by passing the mapped parent contexts to the language model:
\begin{equation}
a = M(q, \{ g(d) \mid d \in C_{\text{small}} \})
\end{equation}

Different chunking and representation methods directly increase the accuracy of RAG pipelines. For example, \citet{lumer2025rethinking} showed that expanding the context to neighboring chunks can increase the win rate to 65\% compared to using single independent chunks. In addition, the adaptive chunking method, where chunks are generated dynamically based on document structure, increases downstream answer correctness by 8--10 percentage points without requiring downstream prompt modifications \cite{demourajunior2026adaptive}. Hierarchical abstract trees like Bridge-RAG \cite{li2026bridgerag} yield up to a 44\% improvement in BLEU scores on complex reasoning tasks by tracing parent-child lineages. Similarly, dynamic reasoning frameworks like LogicRAG \cite{chen2026logicrag} and QPaug \cite{kim2024qpaug} increased string accuracy by 14.7\% absolute (compared to the next best graph baseline) and answer F1 scores by over 34\% relatively.

While small-to-big strategies provide substantial benefits to answer correctness and retrieval recall, they introduce operational and computational overhead. At query time, asynchronous parallel search for neighboring chunks \cite{lumer2025rethinking} keeps latency minimal, adding a negligible $\approx 0.17$\,\text{s} overhead. However, more complex reasoning structures like LogicRAG \cite{chen2026logicrag} more than double the query latency
compared to naive flat retrieval, increasing it from 4.28\,\text{s} to 9.83\,\text{s}. Furthermore, advanced chunking and indexing methods are highly time-consuming during the ingestion phase. The slowest model-driven and layout-aware chunking methods can run for tens of hours, whereas fixed-size chunking completes in under a second \cite{smigielski2026chunking}. In hierarchical structures, generating node summaries and index trees also increases offline preprocessing costs (e.g., scaling from \$1.19 to \$6.12 per average filing in financial domains) \cite{lumer2025rethinking}. Nonetheless, this preprocessing latency is typically acceptable in enterprise environments, as chunking and indexing are offline processes that occur only during initial document ingestion.

\subsection{Semantic Re-ranking}
\label{lit:re_rank}
In RAG pipelines, the chunks might look mathematically relevant (e.g. having high cosine similarity), but lack semantic relevance for the question in hand. The worst case result for this misrelevance is that the answer is wrong and in addition, context is fed with irrelevant data which consumes extra tokens. Semantic re-ranking is an optimization method which uses a multi-stage retrieval where the pipeline first retrieves a large number of chunk candidates using faster and cheaper retrieval models to optimize the recall phase and finally reranks them using reranker model to refine the top most relevant chunks before the generation of the final answer\citep{?, ?}.

% 1. Introductory Paragraph (Concept & Mechanism)
%    - Define semantic reranking as a multi-stage retrieval technique designed to refine query-relevant document ordering.
%    - Explain the core issue: first-stage retrieval (dense vector embeddings or sparse BM25) optimizes for recall, but results in noise and context dilution.
%    - Describe the solution: multi-stage pipelines where a fast first-stage retriever fetches candidates (e.g., top-50/100), and a powerful, computationally intensive cross-encoder or neural reranker refines the top results before generation \citep{moreira2024enhancing, cheng2026enhancing}.
%
% 2. Latency Impact (RQ1 - Latency Profile & Overhead)
%    - First-stage query embedding latency: 5.1\,ms for NV-EmbedQA-E5-v5 (335M parameters) and 19.8\,ms for Mistral-7B-v2 (7.2B parameters) \citep{moreira2024enhancing}.
%    - First-stage bi-encoder retrieval: DPR bi-encoder query embedding is fast (~20\,ms) and FAISS-based retrieval completes in milliseconds (995.0 QPS per thread) \citep{karpukhin2020dense}.
%    - Second-stage reranking latency: NV-RerankQA-Mistral-4B-v3 adds ~266\,ms of latency on average to score 64 candidate passages of 512 tokens \citep{moreira2024enhancing}.
%    - Contrast with agentic RAG: single-pass reranking models like InfoGain-RAG avoid the multiple, slow iterative LLM inference calls of systems like Self-RAG and CRAG, which double query latency from 4.28\,s to 9.83\,s \citep{wang2025infogain, cheng2026enhancing}.
%
% 3. Cost & Token Usage (RQ2 - Operational Cost & Token Reduction)
%    - Model parameter pruning: Pruning a Mistral 7B decoder backbone to 16 layers (creating a 4B parameter model) yields a 50% parameter reduction, reducing GPU memory and serving cost by half, with only a 1% NDCG@10 accuracy penalty \citep{moreira2024enhancing}.
%    - Token compression via context filtering: InfoGain-RAG uses Document Information Gain (DIG) scores to filter out low-value documents before the generation stage, directly reducing the input token size and cost \citep{wang2025infogain}.
%    - Stop-threshold strategies: In financial RAG, document ingestion is stopped once the cumulative probability mass of the cross-encoder reaches 55% (retaining less than 45% of retrieved text), saving token usage by omitting the remaining low-value chunks \citep{cheng2026enhancing}.
%
% 4. Quality & Accuracy Metrics (RQ3 - Retrieval & Generation Quality)
%    - Normalized Discounted Cumulative Gain (NDCG@10): NV-RerankQA-Mistral-4B-v3 achieves a 14% absolute increase in average NDCG@10 across BEIR Q&A datasets (NQ, HotpotQA, FiQA) compared to the second-best reranker (BGE-Reranker-v2-M3) \citep{moreira2024enhancing}.
%    - Exact Match (EM) accuracy: InfoGain-RAG improves Qwen2.5-72B's accuracy on NaturalQA by 17.8% and PopQA by 17.4% over naive RAG. On TriviaQA, LLaMA2-13B reaches 76.2% accuracy, outperforming Self-RAG (69.3%) and CRAG (74.5%) \citep{wang2025infogain}.
%    - Correctness scoring: On the FinDER financial benchmark, cross-encoder reranking raises the correctness rate (scores >= 8 out of 10) from 33.5% to 49.0% (a 15.5 percentage point absolute or 46.3% relative improvement) and reduces completely incorrect answers (score = 1) from 35.3% to 22.5% \citep{cheng2026enhancing}.
%    - Enterprise-wide impact: Systematic survey of 63 studies shows that RAG-augmented enterprise systems achieve up to 87% EM on QA tasks, reducing manual editing time by 30%–50% in field trials \citep{karakurt2025retrieval}.
%
% 5. The Trade-off & Pareto Frontier (RQ4 / Stacking)
%    - Analyze the efficiency-quality equilibrium: reranking shifts the Pareto frontier by providing large quality improvements (+14% to +17.8% accuracy) at the cost of modest execution latency (~266\,ms).
%    - Detail the interaction when stacked: combining hybrid sparse-dense retrieval (e.g., SQLite FTS5 + FAISS vector) with neural reranking delivers optimal results for complex domains (like S&P 500 10-K filings) \citep{cheng2026enhancing}.
%    - Evaluate enterprise compatibility: highlight the importance of commercially licensed models (avoiding research-only MS-MARCO) and the high portability of cross-encoders (onnx exportable) to Java/Spring AI stacks.
%
% =========================================================================

Semantic re-ranking is a multi-stage retrieval technique designed to refine query-relevant document ordering by balancing retrieval recall and precision. In this architecture, a fast first-stage retriever retrieves a relatively large candidate set $C_1$ (e.g., $K_1 = 100$ candidate chunks) from the database using bi-encoder dense embeddings or sparse BM25 retrieval:
\begin{equation}
C_1 = R_{K_1}(q, D)
\end{equation}
Subsequently, a more computationally expensive cross-encoder reranking model, $SR(q, d)$, computes the semantic relevance between the query $q$ and each candidate chunk $d \in C_1$. The top-$K_2$ highest scoring candidate chunks ($K_2 \ll K_1$) are selected to form the generation context $C_2$:
\begin{equation}
C_2 = \operatorname{arg\,top-K_2}_{d \in C_1} SR(q, d)
\end{equation}
The language model then generates the final response using the re-ranked and pruned context:
\begin{equation}
a = M(q, C_2)
\end{equation}


\subsection{Prompt Caching}
\label{lit:caching}

Formally, prompt caching can be implemented as either context-level KV-caching or semantic completion caching.

In context-level caching, the constructed prompt $P(q, C)$ is partitioned into a static prefix $P_{\text{static}}$ (e.g., system instructions, few-shot examples, or pre-retrieved context) and a dynamic query-specific suffix $P_{\text{dynamic}}(q)$:
\begin{equation}
P(q, C) = [P_{\text{static}}, P_{\text{dynamic}}(q)]
\end{equation}
When $P_{\text{static}}$ has been pre-computed, its key-value attention states are retrieved from the cache $KV(P_{\text{static}})$, allowing the model to only compute generation over the suffix:
\begin{equation}
a = M(P_{\text{dynamic}}(q) \mid KV(P_{\text{static}}))
\end{equation}

In semantic caching, a database of past queries and answers $Cache = \{(q_i, a_i)\}$ is maintained. The final answer is determined by:
\begin{equation}
a = \begin{cases} a_i & \text{if } \exists (q_i, a_i) \in Cache \text{ s.t. } \operatorname{sim}(E(q), E(q_i)) \ge \tau \\ M(q, R_k(q, D)) & \text{otherwise} \end{cases}
\end{equation}
where $\tau$ is a similarity threshold.


